\section{Conclusion}
This release turns the project from a paper scaffold into a reproducible prototype benchmark for proactive emotional care timing. The repository now includes a tri-modal synthetic emotion benchmark, a tri-modal care-relation benchmark, a timing synthetic benchmark, a log-derived weak-label dataset, a temporal descriptor backbone, structured and neural baselines, a joint timing-to-strategy branch, and auto-generated experiment tables for the paper.

The main empirical conclusion is cautious but useful: the new tri-modal encoder shows that face, motion, and audio fusion can support fine-grained synthetic emotion understanding very strongly, and the care-relation model shows that the same multimodal stack can learn a synthetic notion of care need and care timing. But weaklabel generalization remains the real bottleneck. The best held-out weaklabel timing Macro-F1 is still only 36.29\%, and the best weaklabel strategy Macro-F1 is 22.50\%. The next step is therefore not to overstate the current system, but to close the remaining gaps: raw-video multimodal encoding, personality ablations, stronger weaklabel reconstruction, and a real-pilot evaluation with human feedback.
